\section{Experimental designs not executed}

\label{sec:agenda-prospective-evaluation}

The remaining questions require distinct targets, rather than one aggregate competence score. The following arms are specified in a machine-readable protocol and have \emph{not} been run. The offline runner has no provider adapter, rejects nonzero-spend or network-enabled configurations, and blocks socket connections. The purpose is to make the proposed research inspectable before any new model expenditure or task-condition change.

\paragraph{Local strategic execution and economic value.}
A local task should first declare its endpoint, such as paying a current debt or reaching a specified construction allocation. Related states and source games should be assigned together to development or held-out sets. Engine witnesses can certify feasibility; they cannot certify that the target maximizes eventual winning probability. An action-value study additionally needs a utility, horizon and opponent distribution, with repeated stochastic continuations. We would compare endpoint attainment, encoding failures and execution separately from independently scored explanation content. General regret or exploitability claims require a defined comparator and opponent class. The existing historical states are public and heavily inspected, so evaluation on them alone cannot establish unfamiliar-task transfer.

\paragraph{Adaptation and memory.}
An adaptation experiment should randomize relevant feedback against a matched control and then present comparable later opportunities, followed by a held-out opponent-policy shift. The unit is a complete task sequence, not an individual correction. The treatment effect is improvement after feedback relative to control, with state-driven reactions, invalid outputs and missing opportunities retained. Cross-game adaptation requires an explicit memory-carryover condition. Weight learning requires training. Neither follows from repeated games or successful corrective retries. Feedback and memory manipulations define separate, versioned experimental conditions; they are not changes to the historical prompt.

\paragraph{Opponent transfer and chance.}
Cross-play should vary opponent families and independently sampled seeds while rotating seats within seed blocks. Opponents must first pass the quality checks above and include meaningful negotiation and development behavior. Complete-system performance and a base-model comparison are different estimands; response policies should be controlled or deliberately factorialized. A historical seed does not guarantee matched opportunities once actions alter random-draw consumption. We would report wins, censoring, rule/encoding failures and task-level endpoints, then assess out-of-sample ordering on held-out opponents. Two chosen pilot seeds do not determine a sufficient sample size or justify a confidence interval for the LLM cohort.

\paragraph{General capability and multi-agent capability.}
Question B requires independent measurements of general ability across multiple tiers within several families. Model configuration, inference resources and response policy must be recorded or controlled. The outcome is a prespecified multi-agent battery, distinct from honesty or instruction adherence. A predictive relationship should be tested on held-out models or families and compared with a family/resource-only baseline. Adding four external leaderboard scores to the existing four configurations would not identify the requested relationship. Precision and expected effect size, rather than an arbitrary game-count target, must determine the prospective sampling plan.

\paragraph{Framing and coordination.}
An optional message experiment can randomize truthful complete disclosure, omission and neutral wording while holding economic terms and other delivered facts fixed. Factual endorsement, action choice and subsequent payoff are separate outcomes. Treatment effects would concern the specified manipulation, not a latent intent label. Changes to public messages are new task conditions, not retroactive changes to historical interactions. A communication intervention can test a specified collective mechanism; parallel actions under shared incentives are not by themselves proof that communication caused a coalition.

The historical workload is not a current cost forecast. No new model query has been made for these arms. Learning, unfamiliar-opponent generalization and the general-to-multi-agent capability relationship remain unmeasured until the corresponding designs are executed.
